\section{Global saturation in degree \(113\)}
\label{sec:global-cm-113}

Put \(p=113\), \(h=56\), and choose \(s=[3]\).
The order of \(2\) modulo \(113\) is \(28\), with \(2=3^{12}\).
Thus
\(\mathcal O=\mathbb Z[\zeta_{113}]\otimes\mathbb Z_2\)
is a product of four unramified degree-\(28\) integer rings.
Complex conjugation \(c=s^{56}\) acts within each factor
as its degree-\(14\) Frobenius. Write
\[
 T=\{u\in1+2\mathcal O:u\,c(u)=1\},\qquad P=(s-1)T.
\]
Then \(T\) has rank \(56\) and torsion \((C_2)^4\).
The generator \(s\) permutes the four signs cyclically.
Since \(s-1\) is invertible on the rational anti-invariant
space, a torsion element in its image must come from a
torsion element. Consequently
\begin{equation}
 P\simeq\mathbb Z_2^{56}\times(C_2)^3.
 \label{eq:113-full-target}
\end{equation}

The integral reflex operator in this primitive coordinate is
\[
\begin{split}
 Q={}&-s^2+s^5-s^6+s^{13}-s^{15}+s^{16}-s^{19}
       +s^{21}-s^{22}+s^{23}\\
    &-s^{24}+s^{26}-s^{30}+s^{37}-s^{39}
       +s^{45}-s^{47}+s^{52}-s^{53}.
\end{split}
\]
In particular \(Q(1)=-1\). On the four sign factors its
reduction is
\begin{equation}
 \overline Q=1+s^2+s^3\pmod{s^4-1},\qquad
 (1+s^2+s^3)^{-1}=1+s+s^2.
 \label{eq:113-sign-inverse}
\end{equation}
All signs therefore belong to the local image.
On the signed cyclic lattice \(\mathbb Z_2[s]/(s^{56}+1)\),
exact fraction-free elimination gives
\[
 \det Q=1028233560187751537651440165000,\qquad
 v_2(\det Q)=3.
\]
Also the Euclidean algorithm over \(\mathbb F_2[s]\) gives
\begin{equation}
 \gcd(\overline Q,s^{56}+1)=s^3+s^2+1.
 \label{eq:113-residual-polynomial}
\end{equation}
The sign calculation and the torsion/free index formula imply
\[
 [T:QT]=[P:QP]=8.
\]
The natural first-layer map is surjective and its target has
dimension three by \eqref{eq:113-residual-polynomial};
hence it is an isomorphism. Inclusion identifies \(P/QP\)
with \(T/QT\), because \(Q\) is multiplication by \(-1\)
on \(T/P\). Thus the quotient is the additive group
\begin{equation}
 \mathscr C=\mathbb F_2[\xi]/(\xi^3+\xi^2+1),
 \qquad s\text{ acts by multiplication by }\xi.
 \label{eq:113-residual-field}
\end{equation}

For the Frobenius calculation it is convenient to use
\[
 L(u)=\log(u)/2\bmod2\mathcal O
     =a+a^2\quad\text{if }u=1+2a.
\]
On \(T\) this is \((1+[2])\theta=(1+s^{12})\theta\).
Since \(\xi\) has order seven, \(1+\xi^{12}\ne0\).
Consequently \(L\), followed by the cyclic quotient
\[
 \beta_k=\overline\zeta_{113}^{\,3^k}
            +\overline\zeta_{113}^{-3^k}\longmapsto\xi^k,
 \qquad0\leq k<56,
\]
also gives an isomorphism \(P/QP\to\mathscr C\).
This logarithmic detector does not remove an unresolved
sign quotient: \eqref{eq:113-sign-inverse} has already
included every finite sign.

\begin{theorem}
\label{thm:113-global-saturation}
After any collection of the other-prime scalar kernels,
the degree-\(113\) images of \(R_{\rm quad}\), \(V_{\rm quad}\),
and their arithmetic intersections are all \(P\).
In particular degree \(113\) can be added to the previously
normalized primes \(29\) and \(\Pi\), preserving the full
relative quotient on the \(R\)-group, its maximal pro-\(2\)
quotient on the \(V\)-group, and the original quadratic
restriction kernel \(\mathcal S\) and full relative image.
The new scalar kernel has infinite index on both groups.
\end{theorem}

\begin{proof}
The prime \(227=2\cdot113+1\) splits completely in
\(\mathbb Q(\zeta_{113})\). The residue \(2\) is primitive
modulo \(227\). Choose \(\chi:\mathbb F_{227}^{\times}
\to\mu_{113}\) by \(\chi(2)=\zeta_{113}\), and extend it
by \(\chi(0)=0\). Put
\[
 J=-\sum_{x\in\mathbb F_{227}}\chi(x)\chi(1-x).
\]
The positive residue-character convention corresponds to
the prime at which \(\zeta_{113}\) reduces to \(4\).
The character-by-character Frobenius comparison
\cite[Theorem~8.1, Remark~8.3, Corollary~8.4]{Otsubo2016}
identifies this Jacobi sum with the degree-\(113\) Tate scalar,
with all four dyadic factors labelled compatibly.
In the inverse-character convention of
\cite[definition on p.~201; Lemma~5.1 and its proof]{AnglesNuccio},
the corresponding residue is \(4^{-1}=57\); this changes
the common label by conjugation and does not change the
image argument.

For an explicit integral certificate, write
\(J=-\sum_{j=0}^{112}n_j\zeta_{113}^j\). The counts
\(n_j=\#\{x\ne0,1:\chi(x(1-x))=\zeta_{113}^j\}\)
are obtained by listing the powers of \(2\) in
\(\mathbb F_{227}^{\times}\). Their sum is \(225\).
Their cyclic autocorrelation has constant coefficient
\(673\) and every other coefficient \(446\), giving
\(Jc(J)=673-446=227\).
The unique odd \(n_j\) is \(n_{111}=1\).
Thus
\[
 \delta=\zeta_{113}^{\,2}J\in1+2\mathcal O.
\]
Its norm is the rational unit \(227\). For a rational
\(2\)-adic unit \(v\), \(\log(v)/2\equiv0\pmod2\);
this follows directly from \(a+a^2=0\) for
\(a\in\mathbb F_2\). Therefore \(L(\delta)\) is
\(c\)-invariant, although the raw first layer of \(\delta\)
need not be.

In the cyclic basis \(\beta_k\), the support of
\(L(\delta)\) is the following subset of \(\{0,\ldots,55\}\):
\[
\begin{split}
 \mathcal I=\{&
 0,1,2,5,7,12,14,15,16,17,18,20,22,24,\\
 &26,28,33,34,38,39,40,41,43,45,46,49,50,55\}.
\end{split}
\]
Reduction by \(\xi^3+\xi^2+1\) consequently gives
\begin{equation}
 \pi L(\delta)=\sum_{k\in\mathcal I}\xi^k
             =1+\xi+\xi^2.
 \label{eq:113-jacobi-detector}
\end{equation}
This certificate uses only arithmetic in the two finite
rings \(\mathbb F_{227}\) and
\(\mathbb F_2[\xi]/(\xi^3+\xi^2+1)\).

Choose \(a\) in \(113\)-inertia inducing \(s\), and let \(b\)
be a Frobenius lift at the indicated prime over \(227\).
Set
\[
 z=[a,b]^{e_2},\qquad w=[a,z].
\]
The established mixed-inertia construction places \(w\)
in \(V_{\rm quad}\cap H\); it also gives \(h_q(w)=1\)
for every \(q\ne113\). These assertions are simultaneous.
For example, at another prime degree the two elements
in the outer commutator have linear rank-one scalar
actions, which commute. The construction also kills the
earlier based-word, radical, cubic, and invariant
quadratic-character coordinates.

In degree \(113\), the inner commutator applies \(s-1\)
to the Frobenius scalar, the pro-\(2\) power removes its
odd part, and the outer commutator applies \(s-1\) again.
Equation \eqref{eq:113-jacobi-detector} yields
\[
 \pi L(h_{113}(w))
    =(\xi-1)^2(1+\xi+\xi^2)=1.
\]
The corresponding values on \(awa^{-1}\) and
\(a^2wa^{-2}\) are \(\xi\) and \(\xi^2\).
These three elements span the residual quotient.
The local norm axes already give \(QP\), including the
signs in \eqref{eq:113-sign-inverse}. Multiplication
therefore gives every element of \(P\), by arithmetic
elements with all other scalar coordinates equal to \(1\).
The four image equalities follow.

On the preceding normalized groups \(R_0,V_0\), write
\(R_1=R_0\cap\ker h_{113}\) and
\(V_1=V_0\cap\ker h_{113}\). Matching a scalar value
by an element of \(V_0\cap H\) gives
\[
 R_0=R_1(V_0\cap H),\qquad
 R_1/V_1\simeq R_0/V_0,\qquad G=R_1H.
\]
The same matching preserves the arithmetic image in the
pro-odd quotient. Extension of invariant relative
characters along \(G=R_1H\), and coprime extension along
\(R_1/V_1\), give the stated relative quotient and original
restriction consequences. Finally
\eqref{eq:113-full-target} has positive rank, so the
new kernels have infinite index.
\end{proof}
